\documentclass{article}

% Keep the package name that matches your local ICLR template files.
% If your template folder provides iclr2027_conference.sty, change 2026 to 2027.
\usepackage{iclr2026_conference,times}

\IfFileExists{math_commands.tex}{\input{math_commands.tex}}{}

\usepackage{booktabs}
\usepackage{graphicx}
\usepackage{hyperref}
\usepackage{tikz}
\usetikzlibrary{arrows.meta,calc,positioning,fit,backgrounds}
\usepackage{url}

\title{Are Neural Models Necessary for \\ Irregular Time Series Forecasting?}

% Anonymous submission. Replace after acceptance or when preparing a non-anonymous draft.
\author{Anonymous Authors}

\newcommand{\method}{AutoAnchor}
\newcommand{\family}{Anchor family}
\newcommand{\best}[1]{\textbf{#1}}
\newcommand{\second}[1]{\underline{#1}}

%\iclrfinalcopy

\begin{document}

\maketitle

\begin{abstract}
Irregular time-series forecasting is commonly treated as a setting that requires
specialized neural architectures for handling uneven timestamps, missingness,
and variable-specific observation patterns. This paper asks a deliberately
simple question: are neural models necessary for the standard irregular
forecasting benchmarks? We introduce an \family{} of deterministic, training-free
history anchors and a unified selector, \method{}, evaluated under the same
splits, scaling, masks, and metrics as the APN benchmark protocol. Early results
show that \method{} matches or exceeds APN-style reported performance on
PhysioNet, USHCN, and HumanActivity, while requiring no backpropagation and no
GPU training. The goal of the paper is not to claim that neural models are never
useful, but to show that current benchmark conclusions are incomplete unless
strong non-neural anchor baselines are included.
\end{abstract}

\section{Introduction}

\begin{itemize}
    \item Motivate irregular time-series forecasting (IMTS): clinical records,
    climate series, human activity traces, and other domains where observations
    arrive at uneven times.
    \item Present the dominant assumption: because the data are irregular,
    sparse, and partially observed, strong forecasting requires specialized
    neural architectures.
    \item Introduce the central question of the paper: \emph{how much of current
    IMTS benchmark performance is explained by local history anchoring rather
    than learned temporal representation?}
    \item Position APN as the current strong reference point: a recent, simple,
    effective neural baseline with public scripts, benchmark datasets, and
    reported results.
    \item State the main thesis: a small deterministic \family{} can match or
    outperform recent neural IMTS models under the same APN protocol.
    \item Contributions:
    \begin{itemize}
        \item Define the \family{}: NaiveAnchor, ExpoAnchor, SparseAnchor,
        ERMAnchor, and \method{}.
        \item Show how \method{} performs model selection over a shared anchor
        library without gradient training or test-label tuning.
        \item Re-evaluate APN benchmark datasets using APN-compatible splits,
        scaling, masks, and global MAE/MSE metrics.
        \item Provide ablations and diagnostic analyses that explain when anchor
        methods succeed and when neural models may still be needed.
    \end{itemize}
\end{itemize}

\section{Background and Motivation}

\subsection{Irregular time-series forecasting}

\begin{itemize}
    \item Define the IMTS setting: each variable has an observation history
    $H_v=\{(t_i,x_i)\}_{i=1}^{n_v}$ with uneven timestamps and missing targets.
    \item Explain why common neural approaches use interpolation, attention,
    message passing, latent ODEs, patching, or time-aware aggregation.
    \item Clarify the benchmark task considered here: predict the next observed
    target values for each variable under APN-style masking and scaling.
\end{itemize}

\subsection{The role of simple baselines}

\begin{itemize}
    \item Connect to the broader lesson from recent time-series work: complex
    architectures can be overstated when simple baselines are weak or missing.
    \item Explain why this question is sharper for IMTS: irregularity often
    motivates architectural complexity, but many horizons may still be governed
    by persistence, smoothing, sparse repetition, or local stability.
    \item State the diagnostic goal: an anchor baseline should be strong enough
    that a neural model must demonstrate learning beyond local history.
\end{itemize}

\subsection{APN benchmark protocol}

\begin{itemize}
    \item Summarize the APN benchmark suite: PhysioNet, USHCN, HumanActivity,
    and MIMIC.
    \item State that this paper uses APN loaders, splits, sequence lengths,
    forecast horizons, scaling, masks, and global scaled MAE/MSE wherever
    available.
    \item Separate APN-reported neural baselines from our single-run anchor
    results.
\end{itemize}

\section{Problem Setup}

\subsection{Forecasting task}

\begin{itemize}
    \item For each variable $v$, the model observes an irregular history
    $H_v$ and predicts the next $K$ target values.
    \item Inputs include both values and timestamps; the anchor methods do not
    discard time.
    \item Forecasts are evaluated only where the APN target mask indicates
    observed ground-truth values.
\end{itemize}

\subsection{Metrics}

\begin{itemize}
    \item Report APN-style global masked scaled MAE and MSE.
    \item Emphasize MSE because large errors are important in clinical and
    scientific forecasting.
    \item In final experiments, include standard deviations when APN exposes
    multiple seeds or splits; otherwise explicitly report fixed-split,
    single-run results.
\end{itemize}

\subsection{Fair comparison protocol}

\begin{itemize}
    \item Anchor selection uses only train and validation targets.
    \item The test split is evaluated once after the algorithm is fixed.
    \item Dataset-specific code is restricted to APN-compatible loading and
    metadata; anchor selection does not branch on dataset or variable names.
\end{itemize}

\section{The Anchor Family}

\IfFileExists{docs/tikz/anchor_family.tex}{%
    % Anchor family overview.
% Requires in the paper preamble:
% \usepackage{tikz}
% \usetikzlibrary{arrows.meta,calc,positioning,fit,backgrounds}

\begin{figure*}[t]
\centering

\definecolor{aaInk}{HTML}{253247}
\definecolor{aaMuted}{HTML}{64748B}
\definecolor{aaBlue}{HTML}{2563EB}
\definecolor{aaCyan}{HTML}{0891B2}
\definecolor{aaGreen}{HTML}{16A34A}
\definecolor{aaOrange}{HTML}{EA580C}
\definecolor{aaRed}{HTML}{DC2626}
\definecolor{aaViolet}{HTML}{7C3AED}
\definecolor{aaGray}{HTML}{CBD5E1}

\begin{tikzpicture}[
  x=1cm,
  y=1cm,
  font=\sffamily,
  >=Stealth,
  box/.style={
    draw=#1,
    fill=white,
    rounded corners=2pt,
    line width=.55pt,
    align=left,
    inner xsep=6pt,
    inner ysep=5pt
  },
  box/.default=aaInk,
  panel/.style={
    box=#1,
    fill=#1!5,
    minimum height=6.25cm
  },
  panel/.default=aaInk,
  mini/.style={
    box=#1,
    fill=#1!8,
    minimum height=.92cm
  },
  mini/.default=aaInk,
  tag/.style={
    draw=#1,
    fill=#1!12,
    rounded corners=1.5pt,
    line width=.45pt,
    font=\scriptsize\sffamily,
    inner xsep=4pt,
    inner ysep=2pt
  },
  tag/.default=aaInk,
  arrow/.style={-{Stealth[length=5pt,width=4pt]}, line width=.75pt, draw=aaInk},
  softarrow/.style={-{Stealth[length=4pt,width=3pt]}, line width=.55pt, draw=#1},
  softarrow/.default=aaInk,
  note/.style={font=\scriptsize\sffamily, text=aaMuted, align=center},
  small/.style={font=\scriptsize\sffamily},
  title/.style={font=\bfseries\small\sffamily, text=aaInk},
  subtitle/.style={font=\scriptsize\sffamily, text=aaMuted}
]

% ---------------------------------------------------------------------
% Module 1: input history
% ---------------------------------------------------------------------
\node[panel=aaBlue, minimum width=3.45cm, anchor=north west] (input) at (0,6.5) {};
\node[title, anchor=west] at ($(input.north west)+(0.22,-0.28)$) {1. Irregular history};
\node[subtitle, anchor=west] at ($(input.north west)+(0.22,-0.58)$)
  {$H_v=\{(t_i,x_i)\}_{i=1}^{n_v}$ per variable};

\begin{scope}[shift={(0.35,4.82)}]
  \draw[aaGray, line width=.4pt] (0,0) -- (2.75,0);
  \draw[aaGray, line width=.4pt] (0,0) -- (0,1.05);
  \foreach \x/\y in {.12/.18,.42/.65,.88/.48,1.34/.92,1.86/.58,2.48/.74} {
    \draw[aaBlue, fill=white, line width=.65pt] (\x,\y) circle (.055);
  }
  \draw[aaBlue, line width=.75pt]
    (.12,.18) -- (.42,.65) -- (.88,.48) -- (1.34,.92) -- (1.86,.58) -- (2.48,.74);
  \foreach \x in {.12,.42,.88,1.34,1.86,2.48} {
    \draw[aaBlue!55, line width=.35pt] (\x,-.06) -- (\x,.02);
  }
  \node[small, text=aaMuted, anchor=west] at (0,-.35) {timestamps are uneven};
\end{scope}

\node[mini=aaBlue, minimum width=2.95cm, text width=2.45cm, anchor=north west] at ($(input.north west)+(0.25,-2.25)$) {
  \textbf{APN protocol}\\[-1pt]
  \scriptsize same split, scaler, mask
};
\node[mini=aaBlue, minimum width=2.95cm, text width=2.45cm, anchor=north west] at ($(input.north west)+(0.25,-3.35)$) {
  \textbf{Forecast request}\\[-1pt]
  \scriptsize next $K$ target values
};
\node[tag=aaBlue, anchor=south west] at ($(input.south west)+(0.25,0.25)$)
  {values + timestamps};

% ---------------------------------------------------------------------
% Module 2: anchor candidate bank
% ---------------------------------------------------------------------
\node[panel=aaGreen, minimum width=5.35cm, anchor=north west] (bank) at (4.05,6.5) {};
\node[title, anchor=west] at ($(bank.north west)+(0.22,-0.28)$) {2. Shared anchor bank};
\node[subtitle, anchor=west] at ($(bank.north west)+(0.22,-0.58)$)
  {candidate forecasts are built without gradient training};

\node[mini=aaGreen, minimum width=4.75cm, text width=4.22cm, anchor=north west] (naive) at ($(bank.north west)+(0.30,-1.08)$) {
  \textbf{NaiveAnchor}\\[-1pt]
  \scriptsize repeat the last observed value
};
\node[mini=aaCyan, minimum width=4.75cm, text width=4.22cm, anchor=north west] (expo) at ($(bank.north west)+(0.30,-2.22)$) {
  \textbf{ExpoAnchor}\\[-1pt]
  \scriptsize exponentially weight recent observations
};
\node[mini=aaOrange, minimum width=4.75cm, text width=4.22cm, anchor=north west] (sparse) at ($(bank.north west)+(0.30,-3.36)$) {
  \textbf{SparseAnchor}\\[-1pt]
  \scriptsize use repetition, mode, or sparse-history persistence
};
\node[mini=aaViolet, minimum width=4.75cm, text width=4.22cm, anchor=north west] (extra) at ($(bank.north west)+(0.30,-4.50)$) {
  \textbf{Additional candidates}\\[-1pt]
  \scriptsize recent means, trimmed means, trend, phase-nearest, mixtures
};

\node[tag=aaGreen, anchor=south west] at ($(bank.south west)+(0.30,0.25)$)
  {one library for all datasets};
\node[tag=aaGreen, anchor=south east] at ($(bank.south east)+(-0.30,0.25)$)
  {CPU friendly};

% ---------------------------------------------------------------------
% Module 3: selection layer
% ---------------------------------------------------------------------
\node[panel=aaViolet, minimum width=5.20cm, anchor=north west] (select) at (10.05,6.5) {};
\node[title, anchor=west] at ($(select.north west)+(0.22,-0.28)$) {3. Selection layer};
\node[subtitle, anchor=west] at ($(select.north west)+(0.22,-0.58)$)
  {choose an anchor before touching the held-out test set};

\node[mini=aaViolet, minimum width=4.58cm, minimum height=1.25cm, text width=4.05cm, anchor=north west] (erm) at ($(select.north west)+(0.32,-1.15)$) {
  \textbf{ERMAnchor}\\[-1pt]
  \scriptsize $a_v^\star=\arg\min_{a\in\mathcal{A}}\mathrm{MSE}_{tr+val}(a)$\\[-1pt]
  \scriptsize optional shrinkage $\beta_v$ is fit only on non-test data
};

\node[mini=aaRed, minimum width=4.58cm, minimum height=1.55cm, text width=4.05cm, anchor=north west] (auto) at ($(select.north west)+(0.32,-3.05)$) {
  \textbf{AutoAnchor}\\[-1pt]
  \scriptsize start from ERMAnchor, then apply generic history-shape priors\\[-1pt]
  \scriptsize density, repetition, phase-like recurrence, local stability
};

\node[box=aaInk, fill=aaInk!4, minimum width=4.58cm, minimum height=.75cm, text width=4.05cm, anchor=north west] (out) at ($(select.north west)+(0.32,-5.05)$) {
  \textbf{Forecast}\\[-1pt]
  \scriptsize selected anchor $\rightarrow$ next $K$ values $\rightarrow$ APN MSE/MAE
};

% ---------------------------------------------------------------------
% Connectors
% ---------------------------------------------------------------------
\draw[arrow] ($(input.east)+(0,.35)$) -- ($(bank.west)+(0,.35)$);
\node[note, fill=white, inner xsep=2pt] at (3.72,4.92) {extract\\history};

\draw[arrow] ($(bank.east)+(0,.30)$) -- ($(select.west)+(0,.30)$);
\node[note, fill=white, inner xsep=2pt] at (9.74,4.92) {score\\candidates};

\draw[softarrow=aaGreen] (naive.east) .. controls +(0.55,0.08) and +(-.55,.45) .. (erm.west);
\draw[softarrow=aaCyan] (expo.east) .. controls +(0.65,0.00) and +(-.55,.25) .. (erm.west);
\draw[softarrow=aaOrange] (sparse.east) .. controls +(0.65,-.10) and +(-.55,.05) .. (erm.west);
\draw[softarrow=aaViolet] (extra.east) .. controls +(0.65,-.20) and +(-.55,-.18) .. (erm.west);

\draw[softarrow=aaViolet] (erm.south) -- ($(erm.south)+(0,-.35)$) -- ($(auto.north)+(0,.03)$);
\draw[softarrow=aaRed] (auto.south) -- (out.north);

% Direct family outputs
\draw[softarrow=aaGreen, dashed] (naive.south west) .. controls +(-.25,-.65) and +(-.80,.15) .. ($(out.west)+(0,.18)$);
\draw[softarrow=aaCyan, dashed] (expo.south west) .. controls +(-.25,-.75) and +(-.80,.05) .. ($(out.west)+(0,.05)$);
\draw[softarrow=aaOrange, dashed] (sparse.south west) .. controls +(-.25,-.65) and +(-.80,-.10) .. ($(out.west)+(0,-.08)$);

% ---------------------------------------------------------------------
% Bottom comparison strip
% ---------------------------------------------------------------------
\node[box=aaInk, fill=aaInk!3, minimum width=15.25cm, minimum height=.95cm, anchor=north west]
  (strip) at (0,0.05) {};
\node[small, text=aaInk, anchor=west, text width=12.8cm] at ($(strip.west)+(0.25,0.28)$) {
  \textbf{Anchor family ladder:}
  NaiveAnchor and ExpoAnchor test persistence and smoothing;
  SparseAnchor tests repeated sparse histories;
  ERMAnchor tests non-neural selection;
  AutoAnchor is the headline unified method.
};
\node[tag=aaInk, anchor=east] at ($(strip.east)+(-0.22,0.28)$)
  {no backprop, no GPU};

\end{tikzpicture}
\caption{
The Anchor family for irregular time-series forecasting.
All methods consume the same irregular history and use the same APN data splits,
scaling, masks, and metrics. Standalone anchors forecast directly from history,
ERMAnchor selects from the shared candidate bank using only train+validation
risk, and AutoAnchor adds generic history-structure priors before the final
held-out test evaluation.
}
\label{fig:anchor-family}
\end{figure*}
%
}{%
    % Anchor family overview.
% Requires in the paper preamble:
% \usepackage{tikz}
% \usetikzlibrary{arrows.meta,calc,positioning,fit,backgrounds}

\begin{figure*}[t]
\centering

\definecolor{aaInk}{HTML}{253247}
\definecolor{aaMuted}{HTML}{64748B}
\definecolor{aaBlue}{HTML}{2563EB}
\definecolor{aaCyan}{HTML}{0891B2}
\definecolor{aaGreen}{HTML}{16A34A}
\definecolor{aaOrange}{HTML}{EA580C}
\definecolor{aaRed}{HTML}{DC2626}
\definecolor{aaViolet}{HTML}{7C3AED}
\definecolor{aaGray}{HTML}{CBD5E1}

\begin{tikzpicture}[
  x=1cm,
  y=1cm,
  font=\sffamily,
  >=Stealth,
  box/.style={
    draw=#1,
    fill=white,
    rounded corners=2pt,
    line width=.55pt,
    align=left,
    inner xsep=6pt,
    inner ysep=5pt
  },
  box/.default=aaInk,
  panel/.style={
    box=#1,
    fill=#1!5,
    minimum height=6.25cm
  },
  panel/.default=aaInk,
  mini/.style={
    box=#1,
    fill=#1!8,
    minimum height=.92cm
  },
  mini/.default=aaInk,
  tag/.style={
    draw=#1,
    fill=#1!12,
    rounded corners=1.5pt,
    line width=.45pt,
    font=\scriptsize\sffamily,
    inner xsep=4pt,
    inner ysep=2pt
  },
  tag/.default=aaInk,
  arrow/.style={-{Stealth[length=5pt,width=4pt]}, line width=.75pt, draw=aaInk},
  softarrow/.style={-{Stealth[length=4pt,width=3pt]}, line width=.55pt, draw=#1},
  softarrow/.default=aaInk,
  note/.style={font=\scriptsize\sffamily, text=aaMuted, align=center},
  small/.style={font=\scriptsize\sffamily},
  title/.style={font=\bfseries\small\sffamily, text=aaInk},
  subtitle/.style={font=\scriptsize\sffamily, text=aaMuted}
]

% ---------------------------------------------------------------------
% Module 1: input history
% ---------------------------------------------------------------------
\node[panel=aaBlue, minimum width=3.45cm, anchor=north west] (input) at (0,6.5) {};
\node[title, anchor=west] at ($(input.north west)+(0.22,-0.28)$) {1. Irregular history};
\node[subtitle, anchor=west] at ($(input.north west)+(0.22,-0.58)$)
  {$H_v=\{(t_i,x_i)\}_{i=1}^{n_v}$ per variable};

\begin{scope}[shift={(0.35,4.82)}]
  \draw[aaGray, line width=.4pt] (0,0) -- (2.75,0);
  \draw[aaGray, line width=.4pt] (0,0) -- (0,1.05);
  \foreach \x/\y in {.12/.18,.42/.65,.88/.48,1.34/.92,1.86/.58,2.48/.74} {
    \draw[aaBlue, fill=white, line width=.65pt] (\x,\y) circle (.055);
  }
  \draw[aaBlue, line width=.75pt]
    (.12,.18) -- (.42,.65) -- (.88,.48) -- (1.34,.92) -- (1.86,.58) -- (2.48,.74);
  \foreach \x in {.12,.42,.88,1.34,1.86,2.48} {
    \draw[aaBlue!55, line width=.35pt] (\x,-.06) -- (\x,.02);
  }
  \node[small, text=aaMuted, anchor=west] at (0,-.35) {timestamps are uneven};
\end{scope}

\node[mini=aaBlue, minimum width=2.95cm, text width=2.45cm, anchor=north west] at ($(input.north west)+(0.25,-2.25)$) {
  \textbf{APN protocol}\\[-1pt]
  \scriptsize same split, scaler, mask
};
\node[mini=aaBlue, minimum width=2.95cm, text width=2.45cm, anchor=north west] at ($(input.north west)+(0.25,-3.35)$) {
  \textbf{Forecast request}\\[-1pt]
  \scriptsize next $K$ target values
};
\node[tag=aaBlue, anchor=south west] at ($(input.south west)+(0.25,0.25)$)
  {values + timestamps};

% ---------------------------------------------------------------------
% Module 2: anchor candidate bank
% ---------------------------------------------------------------------
\node[panel=aaGreen, minimum width=5.35cm, anchor=north west] (bank) at (4.05,6.5) {};
\node[title, anchor=west] at ($(bank.north west)+(0.22,-0.28)$) {2. Shared anchor bank};
\node[subtitle, anchor=west] at ($(bank.north west)+(0.22,-0.58)$)
  {candidate forecasts are built without gradient training};

\node[mini=aaGreen, minimum width=4.75cm, text width=4.22cm, anchor=north west] (naive) at ($(bank.north west)+(0.30,-1.08)$) {
  \textbf{NaiveAnchor}\\[-1pt]
  \scriptsize repeat the last observed value
};
\node[mini=aaCyan, minimum width=4.75cm, text width=4.22cm, anchor=north west] (expo) at ($(bank.north west)+(0.30,-2.22)$) {
  \textbf{ExpoAnchor}\\[-1pt]
  \scriptsize exponentially weight recent observations
};
\node[mini=aaOrange, minimum width=4.75cm, text width=4.22cm, anchor=north west] (sparse) at ($(bank.north west)+(0.30,-3.36)$) {
  \textbf{SparseAnchor}\\[-1pt]
  \scriptsize use repetition, mode, or sparse-history persistence
};
\node[mini=aaViolet, minimum width=4.75cm, text width=4.22cm, anchor=north west] (extra) at ($(bank.north west)+(0.30,-4.50)$) {
  \textbf{Additional candidates}\\[-1pt]
  \scriptsize recent means, trimmed means, trend, phase-nearest, mixtures
};

\node[tag=aaGreen, anchor=south west] at ($(bank.south west)+(0.30,0.25)$)
  {one library for all datasets};
\node[tag=aaGreen, anchor=south east] at ($(bank.south east)+(-0.30,0.25)$)
  {CPU friendly};

% ---------------------------------------------------------------------
% Module 3: selection layer
% ---------------------------------------------------------------------
\node[panel=aaViolet, minimum width=5.20cm, anchor=north west] (select) at (10.05,6.5) {};
\node[title, anchor=west] at ($(select.north west)+(0.22,-0.28)$) {3. Selection layer};
\node[subtitle, anchor=west] at ($(select.north west)+(0.22,-0.58)$)
  {choose an anchor before touching the held-out test set};

\node[mini=aaViolet, minimum width=4.58cm, minimum height=1.25cm, text width=4.05cm, anchor=north west] (erm) at ($(select.north west)+(0.32,-1.15)$) {
  \textbf{ERMAnchor}\\[-1pt]
  \scriptsize $a_v^\star=\arg\min_{a\in\mathcal{A}}\mathrm{MSE}_{tr+val}(a)$\\[-1pt]
  \scriptsize optional shrinkage $\beta_v$ is fit only on non-test data
};

\node[mini=aaRed, minimum width=4.58cm, minimum height=1.55cm, text width=4.05cm, anchor=north west] (auto) at ($(select.north west)+(0.32,-3.05)$) {
  \textbf{AutoAnchor}\\[-1pt]
  \scriptsize start from ERMAnchor, then apply generic history-shape priors\\[-1pt]
  \scriptsize density, repetition, phase-like recurrence, local stability
};

\node[box=aaInk, fill=aaInk!4, minimum width=4.58cm, minimum height=.75cm, text width=4.05cm, anchor=north west] (out) at ($(select.north west)+(0.32,-5.05)$) {
  \textbf{Forecast}\\[-1pt]
  \scriptsize selected anchor $\rightarrow$ next $K$ values $\rightarrow$ APN MSE/MAE
};

% ---------------------------------------------------------------------
% Connectors
% ---------------------------------------------------------------------
\draw[arrow] ($(input.east)+(0,.35)$) -- ($(bank.west)+(0,.35)$);
\node[note, fill=white, inner xsep=2pt] at (3.72,4.92) {extract\\history};

\draw[arrow] ($(bank.east)+(0,.30)$) -- ($(select.west)+(0,.30)$);
\node[note, fill=white, inner xsep=2pt] at (9.74,4.92) {score\\candidates};

\draw[softarrow=aaGreen] (naive.east) .. controls +(0.55,0.08) and +(-.55,.45) .. (erm.west);
\draw[softarrow=aaCyan] (expo.east) .. controls +(0.65,0.00) and +(-.55,.25) .. (erm.west);
\draw[softarrow=aaOrange] (sparse.east) .. controls +(0.65,-.10) and +(-.55,.05) .. (erm.west);
\draw[softarrow=aaViolet] (extra.east) .. controls +(0.65,-.20) and +(-.55,-.18) .. (erm.west);

\draw[softarrow=aaViolet] (erm.south) -- ($(erm.south)+(0,-.35)$) -- ($(auto.north)+(0,.03)$);
\draw[softarrow=aaRed] (auto.south) -- (out.north);

% Direct family outputs
\draw[softarrow=aaGreen, dashed] (naive.south west) .. controls +(-.25,-.65) and +(-.80,.15) .. ($(out.west)+(0,.18)$);
\draw[softarrow=aaCyan, dashed] (expo.south west) .. controls +(-.25,-.75) and +(-.80,.05) .. ($(out.west)+(0,.05)$);
\draw[softarrow=aaOrange, dashed] (sparse.south west) .. controls +(-.25,-.65) and +(-.80,-.10) .. ($(out.west)+(0,-.08)$);

% ---------------------------------------------------------------------
% Bottom comparison strip
% ---------------------------------------------------------------------
\node[box=aaInk, fill=aaInk!3, minimum width=15.25cm, minimum height=.95cm, anchor=north west]
  (strip) at (0,0.05) {};
\node[small, text=aaInk, anchor=west, text width=12.8cm] at ($(strip.west)+(0.25,0.28)$) {
  \textbf{Anchor family ladder:}
  NaiveAnchor and ExpoAnchor test persistence and smoothing;
  SparseAnchor tests repeated sparse histories;
  ERMAnchor tests non-neural selection;
  AutoAnchor is the headline unified method.
};
\node[tag=aaInk, anchor=east] at ($(strip.east)+(-0.22,0.28)$)
  {no backprop, no GPU};

\end{tikzpicture}
\caption{
The Anchor family for irregular time-series forecasting.
All methods consume the same irregular history and use the same APN data splits,
scaling, masks, and metrics. Standalone anchors forecast directly from history,
ERMAnchor selects from the shared candidate bank using only train+validation
risk, and AutoAnchor adds generic history-structure priors before the final
held-out test evaluation.
}
\label{fig:anchor-family}
\end{figure*}
%
}

\subsection{Design principle}

\begin{itemize}
    \item Treat forecasting as selecting a stable reference point from the
    irregular observed history.
    \item Use timestamp-aware candidate summaries rather than learned latent
    states.
    \item Make the baseline transparent: every forecast should be traceable to a
    specific anchor rule and non-test selection criterion.
\end{itemize}

\subsection{Standalone anchors}

\begin{itemize}
    \item \textbf{NaiveAnchor}: repeat the last observed value; tests pure
    persistence.
    \item \textbf{ExpoAnchor}: exponentially smooth recent observations; tests
    local stability under noise.
    \item \textbf{SparseAnchor}: use sparse-history and repeated-value priors;
    tests whether irregular missingness creates mode-like histories.
\end{itemize}

\subsection{Selection anchors}

\begin{itemize}
    \item \textbf{ERMAnchor}: select the candidate anchor that minimizes
    train+validation scaled MSE.
    \item \textbf{\method{}}: start from ERMAnchor, then apply generic
    history-structure priors such as density, repetition, phase-like recurrence,
    and local stability.
    \item Explain that \method{} is unified: no dataset-name or variable-name
    parameter table is used.
\end{itemize}

\subsection{Computational complexity}

\begin{itemize}
    \item No backpropagation, no neural parameter training, and no GPU
    requirement.
    \item Runtime is dominated by candidate scoring on train+validation and
    direct anchor generation on test samples.
    \item Report wall-clock time and memory footprint against APN and other
    neural baselines.
\end{itemize}

\section{Experiments}

\subsection{Datasets}

\begin{itemize}
    \item PhysioNet P12: clinical ICU measurements with many variables and
    sparse observation patterns.
    \item USHCN: climate observations with lower variable count and strong
    temporal persistence.
    \item HumanActivity: activity trajectories with long histories and dense
    short-horizon structure.
    \item MIMIC: credentialed clinical benchmark to be added once local access
    and preprocessing are finalized.
\end{itemize}

\subsection{Main benchmark results}

\begin{table*}[t]
\caption{Current APN benchmark comparison. Lower is better. Bold marks the best
value in each metric column and underline marks the second best. APN and neural
baseline values are taken from the APN benchmark table; Anchor-family values are
single-run APN-style global masked metrics from \texttt{anchor\_results\_family/anchor\_summary.csv}.}
\label{tab:main-results}
\centering
\scriptsize
\resizebox{\textwidth}{!}{
\begin{tabular}{lcccccccc}
\toprule
Method
& \multicolumn{2}{c}{HumanActivity}
& \multicolumn{2}{c}{USHCN}
& \multicolumn{2}{c}{PhysioNet}
& \multicolumn{2}{c}{MIMIC} \\
\cmidrule(lr){2-3}
\cmidrule(lr){4-5}
\cmidrule(lr){6-7}
\cmidrule(lr){8-9}
& MSE & MAE & MSE & MAE & MSE & MAE & MSE & MAE \\
\midrule
PrimeNet & 4.2507$\pm$0.0041 & 1.7018$\pm$0.0011 & 0.4930$\pm$0.0015 & 0.4954$\pm$0.0018 & 0.7953$\pm$0.0000 & 0.6859$\pm$0.0001 & 0.9073$\pm$0.0001 & 0.6614$\pm$0.0001 \\
NeuralFlows & 0.1722$\pm$0.0090 & 0.3150$\pm$0.0094 & 0.2087$\pm$0.0258 & 0.3157$\pm$0.0187 & 0.4056$\pm$0.0033 & 0.4466$\pm$0.0027 & 0.6085$\pm$0.0101 & 0.5306$\pm$0.0066 \\
CRU & 0.1387$\pm$0.0073 & 0.2607$\pm$0.0092 & 0.2168$\pm$0.0162 & 0.3180$\pm$0.0248 & 0.6179$\pm$0.0045 & 0.5778$\pm$0.0031 & 0.5895$\pm$0.0092 & 0.5151$\pm$0.0048 \\
mTAN & 0.0993$\pm$0.0026 & 0.2219$\pm$0.0047 & 0.5561$\pm$0.2020 & 0.5015$\pm$0.0968 & 0.3809$\pm$0.0043 & 0.4291$\pm$0.0035 & 0.9408$\pm$0.1126 & 0.6755$\pm$0.0459 \\
SeFT & 1.3786$\pm$0.0024 & 0.9762$\pm$0.0007 & 0.3345$\pm$0.0022 & 0.4083$\pm$0.0084 & 0.7721$\pm$0.0021 & 0.6760$\pm$0.0029 & 0.9230$\pm$0.0015 & 0.6628$\pm$0.0008 \\
GNeuralFlow & 0.3936$\pm$0.1585 & 0.4541$\pm$0.0841 & 0.2205$\pm$0.0421 & 0.3286$\pm$0.0412 & 0.8207$\pm$0.0310 & 0.6759$\pm$0.0100 & 0.8957$\pm$0.0209 & 0.6450$\pm$0.0072 \\
GRU-D & 0.1893$\pm$0.0627 & 0.3253$\pm$0.0485 & 0.2097$\pm$0.0493 & 0.3045$\pm$0.0305 & 0.3419$\pm$0.0029 & 0.3992$\pm$0.0011 & 0.4759$\pm$0.0100 & 0.4526$\pm$0.0055 \\
Raindrop & 0.0916$\pm$0.0072 & 0.2114$\pm$0.0072 & 0.2035$\pm$0.0336 & 0.3029$\pm$0.0264 & 0.3478$\pm$0.0019 & 0.4044$\pm$0.0020 & 0.6754$\pm$0.1829 & 0.5444$\pm$0.0868 \\
Warpformer & 0.0449$\pm$0.0010 & 0.1228$\pm$0.0018 & 0.1888$\pm$0.0598 & 0.2939$\pm$0.0591 & \second{0.3056$\pm$0.0011} & 0.3661$\pm$0.0016 & \second{0.4302$\pm$0.0035} & \second{0.4025$\pm$0.0014} \\
tPatchGNN & 0.0443$\pm$0.0009 & 0.1247$\pm$0.0031 & 0.1885$\pm$0.0403 & 0.3084$\pm$0.0479 & 0.3133$\pm$0.0053 & 0.3697$\pm$0.0049 & 0.4431$\pm$0.0115 & 0.4077$\pm$0.0088 \\
GraFITi & 0.0437$\pm$0.0005 & 0.1221$\pm$0.0017 & 0.1691$\pm$0.0093 & 0.2777$\pm$0.0248 & 0.3075$\pm$0.0015 & \second{0.3637$\pm$0.0036} & 0.4359$\pm$0.0455 & 0.4142$\pm$0.0297 \\
APN & \second{0.0421$\pm$0.0001} & 0.1159$\pm$0.0006 & \second{0.1590$\pm$0.0137} & \second{0.2611$\pm$0.0167} & 0.3093$\pm$0.0011 & 0.3650$\pm$0.0026 & \best{0.4292$\pm$0.0027} & \best{0.4016$\pm$0.0016} \\
\midrule
NaiveAnchor & 0.0605 & 0.1225 & 0.3640 & 0.3291 & 0.4154 & 0.3981 & -- & -- \\
ExpoAnchor & \best{0.0420} & \second{0.1148} & 0.5923 & 0.5117 & 0.3204 & 0.3716 & -- & -- \\
SparseAnchor & 0.0605 & 0.1225 & 0.2661 & 0.2812 & 0.4154 & 0.3981 & -- & -- \\
ERMAnchor & 0.0428 & \best{0.1131} & 0.1662 & 0.2635 & \best{0.2962} & \best{0.3556} & -- & -- \\
\method{} & \best{0.0420} & \second{0.1148} & \best{0.1565} & \best{0.2187} & \best{0.2962} & \best{0.3556} & -- & -- \\
\bottomrule
\end{tabular}
}
\end{table*}

\subsection{Anchor-family ladder}

\begin{itemize}
    \item Report NaiveAnchor, ExpoAnchor, SparseAnchor, ERMAnchor, and
    \method{} together.
    \item Use the ladder to show that the headline result is not a black-box
    trick: each added mechanism has a measurable role.
    \item Highlight current pattern: persistence alone is strong on some
    variables, ERM improves PhysioNet, and \method{} is strongest on USHCN.
\end{itemize}

\subsection{Efficiency comparison}

\begin{table}[t]
\caption{Efficiency comparison on USHCN. Neural-model GPU memory, parameter
count, training-step time, and inference-step time are from the APN efficiency
comparison. Anchor-family rows are CPU-only timings from our completed
\texttt{anchor\_results\_family} run; they are reported as total
calibration-plus-forecasting wall time and milliseconds per evaluated target
prediction, not neural batch-step inference time.}
\label{tab:efficiency}
\centering
\scriptsize
\begin{tabular}{lcccccc}
\toprule
Method
& GPU Mem. (GB)
& Params (M)
& Train Step (ms)
& Infer. Step (ms)
& CPU Total (s)
& CPU ms / Pred. \\
\midrule
CRU & 0.89 & 31.25 & 1771.00 & 677.00 & -- & -- \\
GraFITi & 0.78 & 989.44 & 34.92 & 21.51 & -- & -- \\
tPatchGNN & 1.09 & 410.71 & 11.35 & 4.22 & -- & -- \\
APN & 0.19 & 1.97 & 3.89 & 1.46 & -- & -- \\
\midrule
NaiveAnchor & 0.00 & 0.00 & N/A & N/A & 0.0681 & 0.1850 \\
ExpoAnchor & 0.00 & 0.00 & N/A & N/A & 0.0700 & 0.1903 \\
SparseAnchor & 0.00 & 0.00 & N/A & N/A & 0.0725 & 0.1969 \\
ERMAnchor & 0.00 & 0.00 & N/A & N/A & 0.0797 & 0.2166 \\
\method{} & 0.00 & 0.00 & N/A & N/A & 0.0806 & 0.2190 \\
\bottomrule
\end{tabular}
\end{table}

\begin{itemize}
    \item Anchor methods have no neural parameters, no gradient-training step,
    and no GPU requirement.
    \item The CPU timing includes both non-test anchor selection and test
    forecasting for USHCN.
    \item The timing columns should be interpreted by execution mode: neural
    models report GPU batch-step times; Anchor methods report CPU
    per-prediction forecasting cost.
\end{itemize}

\section{Analysis}

\subsection{Why do anchors work?}

\begin{itemize}
    \item Quantify persistence: how close are future targets to last observed
    values, recent means, or EMA anchors?
    \item Stratify by observation density and gap length.
    \item Show that many benchmark targets are locally stable enough that
    history anchoring explains much of the performance.
\end{itemize}

\subsection{Where do anchors fail?}

\begin{itemize}
    \item Identify high-error tails and sudden regime shifts.
    \item Analyze PhysioNet variables where clinical changes break persistence.
    \item Use this section to state what neural models should be tested on if
    they claim dynamic representation learning.
\end{itemize}

\subsection{Per-variable and per-dataset diagnostics}

\begin{itemize}
    \item PhysioNet: heatmap of per-variable MAE/MSE and selected anchor rule.
    \item USHCN: analyze sparse/repeated climate histories and seasonal
    structure.
    \item HumanActivity: analyze dense short-horizon trajectories.
\end{itemize}

\section{Ablations}

\begin{itemize}
    \item Remove EMA candidates.
    \item Remove sparse-history candidates.
    \item Remove trend and phase-nearest candidates.
    \item Compare train-only selection, validation-only selection, and
    train+validation selection.
    \item Compare MSE-selected anchors with MAE-selected anchors.
    \item Verify that no test-label tuning or dataset-name-specific parameter
    choices are used.
\end{itemize}

\section{Discussion}

\subsection{Are neural models necessary?}

\begin{itemize}
    \item Give the careful answer: not always under current IMTS benchmark
    protocols.
    \item Explain that neural models may still be necessary for long-horizon
    distribution shift, cross-variable causal dynamics, and true regime changes.
    \item Reframe \method{} as a diagnostic baseline: a neural model should beat
    strong anchors before claiming to learn useful irregular dynamics.
\end{itemize}

\subsection{Implications for benchmarking}

\begin{itemize}
    \item Future IMTS papers should include NaiveAnchor, ExpoAnchor, ERMAnchor,
    and \method{} as default baselines.
    \item Benchmarks should report performance by history density, target gap,
    variable type, and error tail.
    \item New datasets should stress conditions where local anchoring is
    insufficient.
\end{itemize}

\section{Limitations}

\begin{itemize}
    \item Current Anchor-family results are single-run on the APN fixed protocol;
    standard deviations require repeated exposed splits or seeds.
    \item MIMIC is not yet included in the Anchor-family row because it requires
    credentialed access and preprocessing.
    \item Anchor methods are intentionally limited: they do not learn latent
    representations or cross-variable causal structure.
    \item Results may reveal benchmark structure rather than universal
    superiority of non-neural methods.
\end{itemize}

\section{Conclusion}

\begin{itemize}
    \item Summarize the main finding: simple history anchors are unexpectedly
    strong for irregular time-series forecasting.
    \item State the scientific claim: current IMTS benchmark conclusions are
    incomplete without strong non-neural anchors.
    \item Close with the practical value: transparent forecasts, near-zero
    training cost, and a stronger standard for evaluating future neural models.
\end{itemize}

\subsubsection*{Author Contributions}
TODO.

\subsubsection*{Acknowledgments}
TODO.

\bibliography{references}
\bibliographystyle{iclr2026_conference}

\appendix

\section{Additional Experimental Details}

\begin{itemize}
    \item Dataset preprocessing commands.
    \item Exact APN configuration compatibility.
    \item Anchor candidate definitions.
    \item Calibration and selection pseudocode.
\end{itemize}

\section{Additional Results}

\begin{itemize}
    \item Full per-variable PhysioNet table.
    \item Per-dataset calibration logs.
    \item Runtime and memory measurements.
    \item MIMIC results once available.
\end{itemize}

\end{document}
